\documentclass[aps,prb,twocolumn,superscriptaddress,floatfix,longbibliography]{revtex4-2}

\usepackage{amsmath,amssymb,amsfonts,bm}
\usepackage{graphicx}
\usepackage{hyperref}
\usepackage{mathtools}

\begin{document}
\begin{figure}[t]
\includegraphics[width=0.95\columnwidth]{Fig2.pdf}
\caption{Real part of Hall conductivity experimentally measured in various TRSB ferromagnets . FGT: Fe$_3$GeTe$_2$, NMO: Nd$_3$Mo$_2$O$_7$, SRO: SrRuO$_3$.}
\label{fig:Hall}
\end{figure}

We parametrically plot in Fig.~ experimental results for Re$[\sigma_{xy}(\omega)]$ vs. $|\sigma_{xy}(0)|(\frac{\omega}{\hbar\omega})^2$ for a number of magnetic materials. Although there is naturally lots of scatter in the data, Eq.~ again provides an good guide to the rough scale of the dc Hall effect given the high-frequency optical response or vice versa. Again we have plotted this data by setting $\delta$ as the energy where Im$[\sigma_{xy}(\omega)]$ reaches a maximum.

\section{The model's predictive power?}

Finally, having confirmed Eq.~ for a variety of ideal and non-ideal situations for both theory and experiment we can use this model to make a prediction! Using a zero-area-loop Sagnac interferometer and a $\lambda=1550$ nm (0.8 eV) laser, Xin {\it et al.}~ observed a nonzero Kerr rotation of order $\sim 0.8$rad in YBa$_2$Cu$_3$O$_{6+x}$ that onsets at a temperature $T_s$ that tracks the pseudogap temperature $T^*$ as a function of doping and extrapolates toward a putative quantum critical point near optimal doping, suggestive of a genuine broken-symmetry transition rather than a crossover~.

At least two features of these observations sit awkwardly with thinking of the signal as arising from an ordered state with symmetries consistent with a $c$-axis ferromagnet. First, the signal was reported to have the same sign measured from opposite faces of the crystal. Second, it could not be trained by modest fields near $T_c$; its sign appears instead to be fixed by a $\mathcal{T}$-breaking scale already present well above room temperature~. These observations motivated a number of $\mathcal{T}$-preserving interpretation in terms of gyrotropic chiral charge order~, but those proposals were withdrawn as it was noted that Onsager reciprocity forbids a normal-incidence Kerr response in any time-reversal-symmetric medium~. The conclusion that survives is narrower but firmer --- a finite polar Kerr signal requires broken time-reversal symmetry together with the breaking of all vertical mirrors --- while the microscopic order underlying $T^*$ remains unsettled.

Despite these uncertainties (or rather because of them) it makes sense to investigate this physics in different contexts. If we put aside the possibility that the Hall signal has opposite signs on opposite surfaces, our above discussion indicates that a finite Kerr rotation implies a dc Hall effect. Going from the scale of the $\theta_K\sim 0.8$rad YBa$_2$Cu$_3$O$_{6+\delta}$~ and known parameters of the optical response~ in the same frequency range one finds a Hall conductivity of approximately $2\times 10^{-4}$ Ohm$^{-1}\cdot$cm$^{-1}$ at 0.8 eV. If we assume a BCS-like value for the energy scale associated from pseudogap onset temperature e.g $\delta=2\Delta_{PG}=3.5kT^*$, then we predict that the cuprate superconductors should have a dc anomalous Hall effect of order 0.1 Ohm$^{-1}\cdot$cm$^{-1}$. This is a small value, but well within the scale of the anomalous Hall effect observed in various systems~. The challenge may be that this comes with a very small Hall angle $\theta_H\sim 10^{-5}$ due to the relatively low resistivity of cuprates at low temperature. However, Hall angles this small have been measured in Germanium~ and in various magnetic metals like SrRuO$_3$ near their sign change temperature~ and so we encourage a search for this effect. Note that this estimate depends sensitively on the assumed value of $2\Delta_{PG}$ entering Eq.~ squared. A factor of two uncertainty in this scale changes the predicted Hall conductivity by a factor of four. Even allowing for this, the prediction remains within the range of anomalous Hall effects observed in other systems and should be experimentally accessible.

\section{Conclusion}

In this work we have derived a general relation that connects the low-frequency, quasi-dc Hall response to the high-frequency optical Hall response in TRSB states of matter --- ferromagnets, TRSB superconductors, and conductors in a magnetic field. We show that it is obeyed in the most elementary single-mode case of a sharp cyclotron resonance, as well as in theoretical models of ferromagnets and TRSB superconductors. We also show that it agrees well with experimental data on a variety of ferromagnets. Applying the relation to the Hall conductivity inferred from the Kerr effect observed near the pseudogap temperature of the cuprates, we predict that these systems should likewise exhibit an anomalous Hall effect. The estimated size appears to fall within the measurable range, and we encourage an experimental search for it.
\end{document}
